% GENERATED FILE. Edit canonical lessons, then run npm run generate:books.
% canonical-source-hash: fnv1a64:839140758a4e3f29
% canonical-lessons: FA-C297-bozorgtar, FA-C297-mesl, FA-W297-se, FA-C297-se-sal-ast-ke, FA-R297-comparing-and-time

\chapter{Better Than, And Years Ago}
\label{ch:fa-better-than-and-years-ago}
\hlchaptermodality{297}

\begin{chapteropening}
\textbf{By the end of this chapter:} \emph{I can compare two things with -tar az, name the best with -tarin, say what something is like, and say how long something has lasted or how long ago it happened.}

Complete FA-R297-comparing-and-time: compare two places and say how long ago and for how long.
\end{chapteropening}

\section[\texorpdfstring{\fa{بزرگتر} \fa{از} (bozorgtar az)}{بزرگتر از (bozorgtar az)}]{\fa{بزرگتر} \fa{از} (bozorgtar az) — bigger than, and the biggest}

\label{lesson:FA-C297-bozorgtar}



\begin{quote}
Say \textquotedblleft{}big\textquotedblright{}, \textquotedblleft{}small\textquotedblright{} and \textquotedblleft{}better\textquotedblright{}. (\emph{Bozorg}, \emph{kuchek}, \emph{behtar}.)

\emph{Behtar} already holds the pattern. This lesson opens it up for every adjective you know.
\end{quote}

\begin{grammarlens}[title={Grammar Lens: -tar, and az for \textquotedblleft{}than\textquotedblright{}}]
Add \textbf{-tar} to the adjective, and put \textbf{az}, \textquotedblleft{}from\textquotedblright{}, in front of the thing you compare against:

\begin{quote}
\textbf{\fa{این} \fa{اتاق} \fa{از} \fa{آن} \fa{اتاق} \fa{بزرگتر} \fa{است}.} — \emph{in otâq az ân otâq bozorgtar ast} — This room is bigger than that room.
\end{quote}

\begin{quote}
\textbf{\fa{قطار} \fa{از} \fa{اتوبوس} \fa{زودتر} \fa{میرسد}.} — \emph{qatâr az otobus zudtar miresad} — The train arrives earlier than the bus.
\end{quote}

So \emph{kuchektar}, smaller; \emph{zudtar}, earlier; \emph{dirtar}, later. The \emph{az} phrase usually sits before the \emph{-tar} word; \emph{bozorgtar az ân otâq} is heard as well. \emph{Behtar} is the same pattern on \emph{beh}, \textquotedblleft{}good\textquotedblright{}, a word that lives today mostly inside \emph{behtar} and \emph{behtarin}.
\end{grammarlens}

\begin{grammarlens}[title={Grammar Lens: -tarin, in front of the noun}]
For \textquotedblleft{}the most\textquotedblright{}, add \textbf{-tarin} and move the adjective in front of its noun, with no ezafe between them:

\begin{quote}
\textbf{\fa{بهترین} \fa{غذا}} — \emph{behtarin ghazâ} — the best food
\end{quote}

\begin{quote}
\textbf{\fa{این} \fa{بهترین} \fa{رستوران} \fa{شهر} \fa{است}.} — \emph{in behtarin resturân-e shahr ast} — This is the best restaurant in town.
\end{quote}

Every other adjective follows its noun with the ezafe, \emph{ghazâ-ye khub}. The \emph{-tarin} one jumps ahead.
\end{grammarlens}

\begin{cousinweb}
EtymDB 2.1 traces the suffix \emph{-tar} back to an Indo-European suffix that set one of two things against the other. It finds the same suffix at the end of English \textbf{other}: one of two. \emph{-tarin} is \emph{-tar} with a further \emph{-in}.
\end{cousinweb}

\subsection*{Your turn}
Say these aloud:

\begin{itemize}
\raggedright
  \item \emph{in otâq az ân otâq bozorgtar ast}
  \item \textquotedblleft{}The train arrives earlier than the bus\textquotedblright{} — \emph{qatâr az otobus zudtar miresad}
  \item \textquotedblleft{}the best food\textquotedblright{}, then \textquotedblleft{}the biggest city\textquotedblright{} — \emph{behtarin ghazâ}, \emph{bozorgtarin shahr}
  \item a hope with a comparison — \emph{omidvâram ke fardâ behtar bâshad}
\end{itemize}

\begin{tcolorbox}[breakable,colback=teal!4,colframe=teal!35!black,title={Before you move on}]
Which little word says \textquotedblleft{}than\textquotedblright{}? (\emph{\textbf{Az.}}) \textquotedblleft{}Smaller\textquotedblright{}? (\emph{\textbf{Kuchektar.}}) Where does a \emph{-tarin} word go? (\textbf{In front of its noun}, with no ezafe.) Which English word carries the same old suffix as \emph{-tar}? (\textbf{Other.})
\end{tcolorbox}
\section[\texorpdfstring{\fa{مثل} (mesl)}{مثل (mesl)}]{\fa{مثل} (mesl) — like, as}

\label{lesson:FA-C297-mesl}



\begin{quote}
Say \textquotedblleft{}this room is bigger than that room\textquotedblright{}. (\emph{In otâq az ân otâq bozorgtar ast.}) Then \textquotedblleft{}the best food\textquotedblright{}. (\emph{Behtarin ghazâ.})

\emph{-tar} says how two things differ. Describing a place needs a word for how they are the same.
\end{quote}

\subsection*{You'll want to know: \fa{مثل}}
\textbf{\fa{مثل}} — \emph{mesl} — \textquotedblleft{}like, as\textquotedblright{}.

\begin{quote}
\textbf{\fa{این} \fa{شهر} \fa{مثل} \fa{تهران} \fa{است}.} — \emph{in shahr mesl-e Tehrân ast} — This city is like Tehran.
\end{quote}

It takes the ezafe, \emph{mesl-e}, and then the thing it matches: \emph{mesl-e khâne}, like home; \emph{mesl-e ghazâ-ye mâdar}, like mother's cooking.

One use to keep whole: \textbf{\fa{مثل} \fa{همیشه}}, \emph{mesl-e hamishe}, \textquotedblleft{}as usual\textquotedblright{}, word for word \textquotedblleft{}like always\textquotedblright{}:

\begin{quote}
\textbf{\fa{مثل} \fa{همیشه} \fa{دیر} \fa{رسید}.} — \emph{mesl-e hamishe dir resid} — As usual, he arrived late.
\end{quote}

Say it in one beat, \emph{mesl}, with a clear \emph{s}. Its second letter is new: the next lesson puts it in your hand.

\begin{cousinweb}
\emph{mesl} is borrowed from Arabic \emph{mithl}, \textquotedblleft{}like\textquotedblright{} (EtymDB 2.1), on the Arabic root \emph{m-th-l}, usually glossed \textquotedblleft{}to be like\textquotedblright{}. Arabic says its middle letter as the \emph{th} of English \emph{thin}. Persian says it \emph{s} and keeps the Arabic letter in the spelling, the way it kept \textbf{\fa{ذ}} in \emph{ghazâ} and says it \emph{z}.
\end{cousinweb}

\subsection*{Your turn}
Say these aloud:

\begin{itemize}
\raggedright
  \item \emph{mesl}, then \emph{mesl-e}, with the ezafe
  \item \textquotedblleft{}This city is like Tehran\textquotedblright{} — \emph{in shahr mesl-e Tehrân ast}
  \item \textquotedblleft{}As usual, he arrived late\textquotedblright{} — \emph{mesl-e hamishe dir resid}
  \item sameness, then difference — \emph{mesl-e Tehrân ast, vali kuchektar ast}
\end{itemize}

\begin{tcolorbox}[breakable,colback=teal!4,colframe=teal!35!black,title={Before you move on}]
\textquotedblleft{}Like home\textquotedblright{}? (\emph{\textbf{Mesl-e khâne.}}) \textquotedblleft{}As usual\textquotedblright{}? (\emph{\textbf{Mesl-e hamishe.}}) Which says sameness and which says difference? (\emph{\textbf{Mesl}} for sameness, \emph{\textbf{-tar}} for difference.)
\end{tcolorbox}
\section[\texorpdfstring{\fa{ث} (se)}{ث (se)}]{\fa{ث} (se) — the letter inside \fa{مثل}}

\label{lesson:FA-W297-se}



\begin{quote}
Before the new one: write \fa{ت} and \fa{ص} once each, and name them.

Now say \textbf{\fa{مثل}} (\emph{mesl}), \textquotedblleft{}like\textquotedblright{}. You met it a moment ago. This lesson takes one letter out of it and puts it in your hand.
\end{quote}

\subsection*{Script you'll notice: \fa{ث}}
\textbf{\fa{ث}} — \emph{se}, the sound \emph{s}.

It is the middle letter of \textbf{\fa{مثل}}, joined on both sides, \textbf{\fa{ـثـ}}. Its body is the body of \textbf{\fa{ت}}, from chapter fifteen, but it carries \textbf{three dots above} instead of two. In the middle of a word that body shrinks to a single small tooth, and only the dots tell \textbf{\fa{ث}} from \textbf{\fa{ت}}.

You now hold three ways to write \emph{s}: \textbf{\fa{س}}, \textbf{\fa{ص}} and \textbf{\fa{ث}}. As with the four \emph{z} letters, Persian says them alike, and the spelling remembers which came from where. Nearly all words with \textbf{\fa{ث}} came from Arabic.

\subsection*{Writing: \fa{ث}}
Put your pen on \fa{ث} and follow its line. Copy the shape you can see — slowly, and larger than it is printed.

\begin{quote}
This book does not yet tell you \textbf{where to start this letter or which way to travel}. It is not written down here until it can be written down with a source. Copying what is in front of you needs no such source, and it is how the shape gets into your hand in the meantime.
\end{quote}

\subsection*{Your turn}
\begin{itemize}
\raggedright
  \item \emph{Find:} the tooth and three dots of \fa{ث} inside \textbf{\fa{مثل}}, and point at them
  \item \emph{Trace it:} \fa{ث} standing alone three times, saying \emph{s} as you finish each one
  \item \emph{Write it:} \fa{ت}, then \fa{ث}, from memory: two dots, then three
  \item \emph{Write it:} \textbf{\fa{مثل}}, joined
  \item \emph{Read it:} \textbf{\fa{مثل}} aloud once more
\end{itemize}

\begin{tcolorbox}[breakable,colback=teal!4,colframe=teal!35!black,title={Before you move on}]
Which letter is this — \fa{ث}? (\emph{\textbf{se}}, \emph{s}.) Which word did it come from? (\textbf{\fa{مثل}}, \emph{mesl}, \textquotedblleft{}like\textquotedblright{}.) How many dots, and where? (\textbf{Three, above.})
\end{tcolorbox}
\section[\texorpdfstring{\fa{سه} \fa{سال} \fa{است} \fa{که} (se sâl ast ke)}{سه سال است که (se sâl ast ke)}]{\fa{سه} \fa{سال} \fa{است} \fa{که} (se sâl ast ke) — for three years now, and three years ago}

\label{lesson:FA-C297-se-sal-ast-ke}



\begin{quote}
Say \textquotedblleft{}ago\textquotedblright{}, then \textquotedblleft{}a length of time\textquotedblright{}. (\emph{Pish}, \emph{moddat}.) Then \textquotedblleft{}like Tehran\textquotedblright{}. (\emph{Mesl-e Tehrân.})

An experience has a length. Persian measures it two ways, and the verb changes with them.
\end{quote}

\begin{grammarlens}[title={Grammar Lens: ... ast ke, with the present}]
For a stretch that is still going on, name it, add \emph{ast ke}, \textquotedblleft{}it is that\textquotedblright{}, and keep the verb in the present:

\begin{quote}
\textbf{\fa{سه} \fa{سال} \fa{است} \fa{که} \fa{در} \fa{تهران} \fa{زندگی} \fa{میکنم}.} — \emph{se sâl ast ke dar Tehrân zendegi mikonam} — I have lived in Tehran for three years.
\end{quote}

Word for word, \textquotedblleft{}it is three years that I live in Tehran\textquotedblright{}. English switches to \textquotedblleft{}have lived\textquotedblright{}; Persian stays in the present, because the living has not stopped. \emph{Zendegi kardan} is \textquotedblleft{}to live\textquotedblright{}, \emph{zendegi}, \textquotedblleft{}life\textquotedblright{}, on \emph{kardan}.

\begin{quote}
\textbf{\fa{دو} \fa{هفته} \fa{است} \fa{که} \fa{اینجا} \fa{کار} \fa{میکنم}.} — \emph{do hafte ast ke injâ kâr mikonam} — I've been working here for two weeks.
\end{quote}

To ask: \textbf{\fa{چه} \fa{مدت} \fa{است} \fa{که} \fa{اینجا} \fa{کار} \fa{میکنید؟}}, \emph{che moddat ast ke injâ kâr mikonid?}, \textquotedblleft{}how long have you been working here?\textquotedblright{}.
\end{grammarlens}

\begin{grammarlens}[title={Grammar Lens: pish, with the plain past}]
For a point in the past, the time goes first with \emph{pish} after it, and the verb is the plain past:

\begin{quote}
\textbf{\fa{سه} \fa{سال} \fa{پیش} \fa{به} \fa{تهران} \fa{آمدم}.} — \emph{se sâl pish be Tehrân âmadam} — I came to Tehran three years ago.
\end{quote}

And a stretch that is over takes \emph{barâ-ye}, \textquotedblleft{}for\textquotedblright{}, with the plain past: \emph{barâ-ye do hafte be Tehrân raftam}, \textquotedblleft{}I went to Tehran for two weeks\textquotedblright{}.

So the three fit together: \emph{se sâl pish âmadam. se sâl ast ke injâ zendegi mikonam.}
\end{grammarlens}

\subsection*{Your turn}
\begin{itemize}
\raggedright
  \item \emph{Say it:} \emph{se sâl ast ke dar Tehrân zendegi mikonam}
  \item \emph{Say it:} \textquotedblleft{}I came to Tehran three years ago\textquotedblright{} — \emph{se sâl pish be Tehrân âmadam}
  \item \emph{Ask:} \emph{che moddat ast ke injâ kâr mikonid?}
  \item \emph{Say it:} your own — how long ago you started something, and how long it has gone on
\end{itemize}

\begin{tcolorbox}[breakable,colback=teal!4,colframe=teal!35!black,title={Before you move on}]
Why does \emph{se sâl ast ke ... zendegi mikonam} keep the present? (The living \textbf{has not stopped}.) \textquotedblleft{}Three years ago\textquotedblright{}? (\emph{\textbf{Se sâl pish.}}) \textquotedblleft{}How long?\textquotedblright{} (\emph{\textbf{Che moddat?}})
\end{tcolorbox}
\section[(review)]{Review — two places, and how long}

\label{lesson:FA-R297-comparing-and-time}



\begin{quote}
Say \textquotedblleft{}smaller than that room\textquotedblright{}. (\emph{Az ân otâq kuchektar.}) Then \textquotedblleft{}three years ago\textquotedblright{}. (\emph{Se sâl pish.})
\end{quote}

\begin{grammarlens}[title={What you've built}]
Two chapters, and an experience can now be measured. \emph{-tar az} sets one thing against another, \emph{-tarin} crowns one, and \emph{mesl} says two are alike. \emph{Pish} fixes a moment in the past; \emph{... ast ke} with the present measures a stretch that is still running.

Someone who moved, describing it:

\begin{quote}
\emph{do sâl pish be in shahr âmadam. do sâl ast ke injâ zendegi mikonam. in shahr mesl-e Tehrân ast, vali kuchektar ast. ghazâ-ye injâ behtar ast! avval nârâhat budam, vali ba'd khoshhâl shodam.}
\end{quote}

The last line goes back to chapter two hundred and ninety-six: a mood that lasted, \emph{budam}, and a change that arrived, \emph{shodam}.
\end{grammarlens}

\subsection*{Your turn}
Say these aloud:

\begin{itemize}
\raggedright
  \item the mover's five lines, without stopping
  \item the same about a place you know, with \emph{pish}, \emph{ast ke} and \emph{mesl}
  \item the best thing there — \emph{behtarin...}
  \item a hope to finish — \emph{omidvâram ke...} or \emph{dust dâram...}
\end{itemize}

\begin{tcolorbox}[breakable,colback=teal!4,colframe=teal!35!black,title={Before you move on}]
\textquotedblleft{}I've lived here for two years\textquotedblright{}? (\emph{\textbf{Do sâl ast ke injâ zendegi mikonam.}}) Where does \emph{behtarin} go? (\textbf{In front of its noun.}) Which letter in \textbf{\fa{مثل}} says \emph{s}? (\textbf{\fa{ث}}, \emph{se}.) And which one opens \textbf{\fa{غذا}}? (\textbf{\fa{غ}}, \emph{gheyn}.)
\end{tcolorbox}
